\section{Preliminaries}


  The goal of this note is to illustrate the difference between the following notions of positivity for divisors on a projective variety $X$:
  \begin{center}
    \begin{tikzcd}
      \boxed{\text{base point free}} \arrow[d, Rightarrow] \arrow[r, Rightarrow] & \boxed{\text{(linearly/numerically) effective}} \arrow[d, Rightarrow] \\ 
      \boxed{\text{nef}} \arrow[r, Rightarrow] & \boxed{\text{pseudo-effective}}
    \end{tikzcd}
  \end{center}
  

\subsection{Notations and conventions}

  \begin{notation}
    We will use the following notations throughout this note.
    \begin{itemize}
      \item All schemes are assumed to be separated and excellent.
      \item We denote by $\kk$ an arbitrary field, and $\kkk$ its the algebraic closure or an algebraically closed field.
        We do not restrict the characteristic of $\kk$ unless otherwise specified.
      \item A \emph{variety} is a geometrically integral separated scheme of finite type over $\kk$.
      \item A \emph{divisor} will usually indicate a Cartier divisor, unless otherwise specified.
        However, by \emph{prime divisor} we will always mean a prime Weil divisor.
        The group of Cartier (resp. Weil) divisors is denoted by $\Div(X) = \CaDiv(X)$ (resp. $\WDiv(X) = Z_{d-1}(X)$, where $d = \dim X$).
      \item The linear equivalence of divisors (or Weil divisors) will be denoted by $\sim$.
        The group of linear equivalence classes of divisors (resp. Weil divisors) is denoted by $\Pic(X)$ (resp. $\CH_{d-1}(X)$).
        The $\bbQ$-linear equivalence of divisors (or Weil divisors) will be denoted by $\sim_\bbQ$.
        Note that for integral divisors, we have $D \sim_\bbQ D' \iff D - D'$ is torsion in $\Pic(X)$.
       
      \item Suppose that $X$ is projective.
        The \emph{numerical equivalence} of divisors is defined by intersecting with curves, denoted by $\equiv$.
        Since for numerical equivalence, there is no issue of torsion, we will not distinguish between $\bbR$, $\bbQ$ and $\bbZ$-numerical equivalence.
        The group of numerical equivalence classes of divisors is denoted by $\NS(X)$, and $\NS(X)_\bbQ = \NS(X) \otimes_\bbZ \bbQ$, and $\upN^1(X) = \NS(X) \otimes_\bbZ \bbR$.
        The dual space of $\upN^1(X)$ is denoted by $\upN_1(X)$.
        We have a natural surjective map $Z_1(X)_\bbR \to \upN_1(X)$ induced by the intersection pairing.

    \end{itemize}
    
    \Yang{}
  \end{notation}
  
  As real vectors spaces, we have $\dim_\bbR \upN^1(X) = \dim_\bbR \upN_1(X) < \infty$; see \cite[Theorem 2.5.1]{Dang20degrees}.
    Although Dang only deals with the case when $\kk$ is algebraically closed and $X$ is normal, and his definition is slightly different, the same proof works in our setting by basing change to $\kkk$.\Yang{}

  \begin{remark}
    Although I try to work over an arbitrary field $\kk$, 
    one should be careful when $\kk$ is not algebraically closed or when $\kk$ has positive characteristic, since I may not be aware of some subtleties in these cases.
     \Yang{}
  \end{remark}

  \Yang{support of a cartier divisor when $X$ is not normal}




\subsection{Varieties over non-closed fields}

  In this subsection, we review some facts about varieties over non-closed fields.
  Let $\kk$ be a field and $\kkk$ be its algebraic closure.

  Many properties of schemes are not stable under base change to $\kkk$, for example, being integral, irreducible, reduced, connected, etc.

  \begin{example}
    Let $X = \Spec \bbR[x,y]/(x^2 + y^2)$.
    Then $X$ is integral.
    However, $X_\bbC = \Spec \bbC[x,y]/(x^2 + y^2)$ has two irreducible components defined by $y = ix$ and $y = -ix$ respectively, and hence is not integral.
    %  \Yang{}
  \end{example}

  \begin{example}\label{eg:algebraic_extension_and_geometrically_integrality}
    Let $\kk$ be a field and $\kk' = \kk(\alpha) = \kk[x]/(f)$ be a finite algebraic extension, where $f$ is an irreducible polynomial of degree $n \geq 2$ over $\kk$.
    Then $X = \Spec \kk'$ is integral, but $X_\kkk = \Spec \kk' \ten_\kk \kkk$ is not integral since 
    \[ \kk' \ten_\kk \kkk \cong \kkk[x]/(f) \cong \prod_{i=1}^k \kkk[x]/((x - \alpha_i)^{m_i}), \] 
    where $\alpha_i$ are the distinct roots of $f$ in $\kkk$ and $m_i$ are their multiplicities.

    More generally, suppose $k \geq 2$ (i.e., $f$ has at least two distinct roots in $\kkk$) and let $Y$ be any integral scheme over $\kk'$.
    Regard $Y$ as a scheme over $\kk$ via the composition $Y \to \Spec \kk' \to \Spec \kk$.
    Then $Y$ is integral, but $Y \times_\kk \kkk = Y \times_{\kk'} (\kk' \ten_\kk \kkk)$ is not integral since it has at least two irreducible components $Y \times_{\kk'} \kkk[x]/((x - \alpha_i)^{m_i})$ for $i = 1, 2$.
    %  \Yang{}
  \end{example}

  \begin{definition}
    Let $X$ be a scheme of finite type over $\kk$.
    We say $X$ is \emph{geometrically integral} (resp. \emph{geometrically irreducible}, \emph{geometrically reduced}, \emph{geometrically connected}) if $X_{\kk'} = X \times_\kk \kk'$ is integral (resp. irreducible, reduced, connected) for any field extension $\kk'/\kk$.
    One can similarly define other notions of geometric properties, for example, \emph{geometrically normal}, \emph{geometrically regular}, etc.
  \end{definition}

  A fact is that we only need to check the algebraical closure of $\kk$ (\cite[\href{https://stacks.math.columbia.edu/tag/030V}{Tag 030V}, \href{https://stacks.math.columbia.edu/tag/037K}{Tag 037K} and \href{https://stacks.math.columbia.edu/tag/0FWF}{0FWF}]{Stacks}).

  For geometrically irreducibility, one only needs to check the separable closure of $\kk$ in $\fkk(X)$.

  \begin{lemma}[{ref. \cite[\href{https://stacks.math.columbia.edu/tag/0G33}{Tag 0G33} and \href{https://stacks.math.columbia.edu/tag/054Q}{054Q}]{Stacks}}]\label{lem:geometrically_irreducible_and_separably_close_in_function_field}
    Let $X$ be an integral scheme of finite type over $\kk$ and $\kk' \subset \fkk(X)$ be the field consisting of elements of $\fkk(X)$ that are separably algebraic over $\kk$.
    Then $X$ is geometrically irreducible if and only if $\kk' = \kk$.
  \end{lemma}
  % \begin{proof}
  %   By \cite[\href{https://stacks.math.columbia.edu/tag/0G33}{Tag 0G33}]{Stacks}, we have $X_\kkk = X \times_\kk \kkk \cong X \times_{\Spec \kk'} \Spec \kkk$.
  %   Then $X_\kkk$ is irreducible if and only if $X$ is irreducible over $\kk'$, which is equivalent to $\kk' = \kk$ since $X$ is integral over $\kk$.
  %   \Yang{}
  % \end{proof}

  The case of geometrically reduced is more subtle, namely, we cannot only check the algebraic part $\kk'/\kk$ of $\fkk(X)/\kk$; see the following example.

  \begin{example}[{ref. \cite[\href{https://stacks.math.columbia.edu/tag/020F}{Tag 020F}]{Stacks}}]
    Let $\kk = \bbF_p(s,t)$ and $\KK = \Frac (\kk[x,y]/(sx^p + ty^p + 1))$.
    Then $\kk$ is algebraically closed in $\KK$ and $\KK/\kk$ is not separable.
    See \cref{eg:fibration_with_non_reduced_fiber} for a geometric views.
    \Yang{}
  \end{example}

  \begin{example}\label{eg:fibration_with_non_reduced_fiber}
    Suppose $\characteristic \kk = p > 0$.
    Let $Y = \Spec \kk[s,t]$ and $X = \Spec \kk[s,t,x,y]/(sx^p + ty^p - 1)$.
    Since $Y$ is normal and $\KK := \fkk(Y) = \kk(s,t)$ is algebraically closed in $\fkk(X)$, the morphism $f: X \to Y$ is a fibration.

    The generic fiber of $f$ is given by $\Spec \KK[x,y]/(sx^p + ty^p - 1)$.
    Over corresponding geometric point $\Spec \KKK$ of $\Spec \KK$, the fiber is given by $\Spec \KKK[x,y]/(sx^p + ty^p - 1) = \Spec \KKK[x,y]/((s^{1/p}x + t^{1/p}y - 1)^p)$, which is not reduced.
     \Yang{}
  \end{example}

  We list the following results without proof, and refer to \cite[\href{https://stacks.math.columbia.edu/tag/030I}{Tag 030I}, \href{https://stacks.math.columbia.edu/tag/05DS}{Tag 05DS}, \href{https://stacks.math.columbia.edu/tag/05DT}{Tag 05DT}, \href{https://stacks.math.columbia.edu/tag/05DU}{Tag 05DU} and \href{https://stacks.math.columbia.edu/tag/035U}{Tag 035U}]{Stacks} for details.
  
  \begin{definition}\label{def:separable_field_extensions_and_morphisms}
    Let $\KK/\kk$ be a field extension.
    We say $\KK/\kk$ is \emph{separable} if $\KK \ten_\kk \bfL$ is a reduced ring for every field extension $\kk \subseteq \bfL$.

    Let $f: X \to Y$ be a dominant morphism between integral schemes.
    We say $f$ is \emph{separable} if the field extension $\fkk(Y) \subseteq \fkk(X)$ is separable.
  \end{definition}

  \begin{lemma}\label{lem:criterion_of_geometrical_reducedness_by_separability}
    Let $\kk$ be a field, $X$ be an integral scheme of finite type over $\kk$, and $\KK/\kk$ be a finitely generated field extension.
    \begin{enumerate}
      \item If $\fkk(X)/\kk$ is separable, then $X$ is geometrically reduced (ref. \cite[\href{https://stacks.math.columbia.edu/tag/030U}{Tag 030U}]{Stacks}).
      \item If $\characteristic \kk = 0$, then any extension $\KK/\kk$ is separable (ref. \cite[\href{https://stacks.math.columbia.edu/tag/030Z}{Tag 030Z}]{Stacks}).
      \item The extension $\KK/\kk$ is separable if and only if there exists a transcendence basis $\{x_1, \dots, x_r\}$ of $\KK/\kk$ such that $\KK$ is a separable algebraic extension of $\kk(x_1, \dots, x_r)$ (ref. \cite[\href{https://stacks.math.columbia.edu/tag/030X}{Tag 030X}]{Stacks}; \cite[Theorem 26.2]{Mat89}).
      %  \Yang{}
    \end{enumerate}
  \end{lemma}


\subsection{Line bundles on varieties over non-closed fields}



  \begin{proposition}


    We have 
    \begin{enumerate}
      \item $\Pic(X)_\bbQ = \Pic(X_\kkk)_\bbQ$ under $\pi^*: X_\kkk \to X$.
      \item $X$ admit an ample divisor iff $X_\kkk$.
      \item Suppose that $X$ is projective.
        Then under $\pi^*: X_\kkk \to X$, we have $\upN^1(X) = \upN^1(X_\kkk)$, so is $\Psef^1, \Nef^1$.
    \end{enumerate}
    
    
    \Yang{}
  \end{proposition}

  
  When dealing with varieties over non-closed fields, we need to be careful about the degree of a closed point.

  \begin{example}
    Let $X = \Spec \bbR[x, y]/(x^2 + y^2 - 1)$ be the ``unit circle''.
    Note that $X$ is a open subscheme of $\bbP^1_\bbR \cong \Proj \bbR[x,y,z]/(x^2 + y^2 - z^2)$.
    
    We have $\Pic X = \bbZ/2\bbZ$ rather than $0$ since $X = \bbP^1_\bbR \setminus \{P\}$, where $P$ is a closed point with $\rkk(P) = \bbC$ and hence $\calO_{\bbP^1_\bbR}(P) \cong \calO_{\bbP^1_\bbR}(2)$.
    Then $\Pic X \cong \bbZ/2\bbZ$ follows from the localization sequence of Picard groups since $P$ is irreducible closed subscheme of $\bbP^1_\bbR$.
    And we have $\Pic X_\bbC = 0$ since $X_\bbC \cong \bbP^1_\bbC \setminus \{P_1, P_2\}$, where $P_1$ and $P_2$ are two distinct $\bbC$-points mapping to $P$.
    % \Yang{}
  \end{example}







% \subsection[b-divisors]{$\bfb$-divisors}

%   Let $X$ be a projective variety over $\kk$.
%   The partially ordered set of all birational models of $X$ forms a directed system, we denote it by $\frakX$.
%   Note that here $\frakX$ is just a notation for a set (with a partial order), and it has no geometric structure.
%   Also note that $\frakX$ only depends on the field extension $\kk(X)/\kk$.
%   Sometimes we will also consider the sub-directed system of all normal (resp. smooth) birational models of $X$, which we will denote by $\frakX_{\mathrm{nor}}$ (resp. $\frakX_{\sm}$).
%   We define the following datum of $\frakX$.

%   \begin{definition}
%     We define 
%     \[ \Div(\frakX) \coloneqq \CaDiv(\frakX) \coloneqq \dirlim_{Y \in \frakX} \Div(Y). \]
%     Similar definitions for $\Pic(\frakX)$, $\NS(\frakX)$, $\upN^1(\frakX)$, etc.
%     Elements of $\Div(\frakX)$ (resp. $\Pic(\frakX)$, $\NS(\frakX)$, $\upN^1(\frakX)$) are called \emph{$\bfb$-Cartier divisors} (resp. \emph{linearly equivalent classes of $\bfb$-Cartier divisors}, \emph{numerical equivalence classes of $\bfb$-Cartier divisors}, \emph{real numerical equivalence classes of $\bfb$-Cartier divisors}).

%     For cycles, we define
%     \[ Z_{k}(\frakX) \coloneqq \invlim_{Y \in \frakX} Z_{k}(Y), \]
%     Similar definitions for $\CH_{k}(\frakX)$, $\upN_k(\frakX)$, etc.
%     Elements of $Z_{k}(\frakX)$ (resp. $\CH_{k}(\frakX)$, $\upN_1(\frakX)$) are called \emph{$\bfb$-cycles} (resp. \emph{$\bfb$-cycle classes}, \emph{numerical equivalence classes of $\bfb$-cycles}).
%     If we only consider normal birational models, we will also use the notation $\WDiv(\frakX) \coloneqq Z_{d-1}(\frakX_\nor)$, $\WCl(\frakX) \coloneqq \CH_{d-1}(\frakX_\nor)$.
%     %  \Yang{}
%   \end{definition}

%   For the cohomological objects, for instance, $\Div(\frakX)$, we have the following description:
%   an element of $\Div(\frakX)$ is represented by a divisor $D_Y$ on some birational model $Y$ of $X$, and two divisors $D_Y$ and $D_{Y'}$ represent the same element in $\Div(\frakX)$ if there exists a birational model $Z$ dominating both $Y$ and $Y'$ such that the pullbacks of $D_Y$ and $D_{Y'}$ to $Z$ are equal.

%   For the homological objects, for instance, $\WDiv(\frakX)$, we have the following description:
%   an element of $\WDiv(\frakX)$ is represented by a collection of divisors $D_Y$ on all birational models $Y$ of $X$ such that for any birational morphism $f: Y' \to Y$, we have $D_{Y} = f_* D_{Y'}$.




%   % \begin{lemma}
%   %   There is a one-to-one correspondence between the following two sets:
%   %   \begin{enumerate}
%   %     \item transcendental, finitely generated field extensions $\KK/\kk$,
%   %     \item birational equivalence classes of normal projective varieties over $\kk$.
%   %   \end{enumerate}
%   %    \Yang{}
%   % \end{lemma}

%   % \begin{definition}
%   %   Let $\KK/\kk$ be a transcendental, finitely generated field extension.
%   % \end{definition}

%   % \begin{definition}[Zariski-Riemann space]
%   %   Let $\KK/\kk$ be a transcendental, finitely generated field extension.
%   %   The \emph{Zariski-Riemann space} $\frakX$ of $\KK/\kk$ is the projective limit of all projective varieties $X$ over $\kk$ such that $\kk(X) \cong \KK$, i.e.
%   %   \[ \frakX := \varprojlim_{X} X. \]
%   %    \Yang{}
%   % \end{definition}


%   % \begin{definition}
%   %   Let $X$ be a normal variety.
%   %   A \emph{$\bfb$-divisor} $\bfD$ on $X$ is a collection of divisors $\bfD_Y$ on all birational models $Y$ of $X$ such that for any birational morphism $f: Y' \to Y$, we have $\bfD_{Y'} = f^* \bfD_Y$.
%   % \end{definition}